\chapter{Référence : les commandes et le vocabulaire}

Ce chapitre est une table. Il se consulte, il ne se lit pas --- et il est relevé
sur la version 2.4.0.

\section{Les vingt-sept commandes}

\begin{center}
\begin{tabular}{@{}llp{6.6cm}@{}}
\toprule
\textbf{Commande} & \textbf{Alias} & \textbf{Ce qu'elle fait} \\
\midrule
\texttt{build} & \texttt{b} & compile le workspace ou un projet \\
\texttt{run} & \texttt{r} & exécute un projet \\
\texttt{test} & \texttt{t} & lance les tests unitaires \\
\texttt{gdb} & \texttt{g}, \texttt{debug} & lance le débogueur \\
\texttt{clean} & \texttt{c} & supprime objets et artefacts \\
\texttt{rebuild} & --- & \texttt{clean} puis \texttt{build} \\
\texttt{watch} & \texttt{w} & reconstruit à chaque modification \\
\texttt{info} & \texttt{i} & ce que Jenga a compris \\
\texttt{gen} & --- & produit CMake, Visual Studio, Makefile \\
\texttt{workspace} & \texttt{init} & crée un espace de travail \\
\texttt{project} & \texttt{create} & crée un projet \\
\texttt{file} & \texttt{add} & ajoute fichiers ou dépendances \\
\texttt{install} & --- & installe dépendances et chaînes locales \\
\texttt{keygen} & \texttt{k} & génère une clé de signature \\
\texttt{sign} & \texttt{s} & signe un artefact \\
\texttt{docs} & \texttt{d} & la documentation \\
\texttt{package} & --- & construit puis empaquette \\
\texttt{deploy} & --- & installe sur une cible \\
\texttt{publish} & --- & envoie vers une distribution \\
\texttt{profile} & --- & profilage \\
\texttt{bench} & --- & mesures \\
\texttt{config} & --- & configuration de Jenga \\
\texttt{examples} & \texttt{e} & les exemples fournis \\
\texttt{ide-setup} & \texttt{ide} & configure l'éditeur \\
\texttt{compile-flags} & --- & les drapeaux réels, en JSON \\
\texttt{help} & \texttt{h} & la liste \\
\bottomrule
\end{tabular}
\end{center}

\begin{nkretenir}
\textbf{Les alias ne sont pas tous d'une lettre}, et deux d'entre eux sont des
\emph{synonymes} plutôt que des raccourcis : \texttt{debug} pour \texttt{gdb},
\texttt{ide} pour \texttt{ide-setup}.

Et deux paires sont réversibles : \texttt{workspace}/\texttt{init},
\texttt{project}/\texttt{create}, \texttt{file}/\texttt{add}. Employez le même
dans une équipe : deux orthographes rendent une recherche dans les scripts
incomplète.
\end{nkretenir}

\section{Le vocabulaire, par famille}

\subsection*{Structure}

\begin{center}
\begin{tabular}{@{}ll@{}}
\toprule
\texttt{workspace(nom)} & le bloc racine \\
\texttt{project(nom)} & un artefact \\
\texttt{filter(expr)} & une condition \\
\texttt{include(chemin)} & exécute un autre \texttt{.jenga} en ligne \\
\texttt{unitest()} & configure le cadre de test \\
\texttt{toolchain(\dots)} & déclare une chaîne d'outils \\
\bottomrule
\end{tabular}
\end{center}

\subsection*{Genre et langage}

\begin{center}
\begin{tabular}{@{}ll@{}}
\toprule
\texttt{consoleapp()} \texttt{windowedapp()} & exécutables \\
\texttt{staticlib()} \texttt{sharedlib()} & bibliothèques \\
\texttt{testsuite()} & suite de tests \\
\texttt{kind(k)} & la forme générique \\
\texttt{language(l)} & \texttt{"C"} ou \texttt{"C++"} \\
\texttt{cppdialect} \texttt{cdialect} & la norme \\
\bottomrule
\end{tabular}
\end{center}

\subsection*{Fichiers et chemins}

\begin{center}
\begin{tabular}{@{}ll@{}}
\toprule
\texttt{location} & où vit le projet \\
\texttt{files} \texttt{excludefiles} \texttt{excludemainfiles} & les sources \\
\texttt{includedirs} \texttt{externalincludedirs} \texttt{sysincludedirs} & en-têtes \\
\texttt{libdirs} \texttt{syslibdirs} & où chercher les bibliothèques \\
\texttt{objdir} \texttt{targetdir} \texttt{bindir} \texttt{targetname} & les sorties \\
\texttt{dependfiles} \texttt{embedresources} & fichiers copiés / intégrés \\
\bottomrule
\end{tabular}
\end{center}

\subsection*{Dépendances et drapeaux}

\begin{center}
\begin{tabular}{@{}ll@{}}
\toprule
\texttt{dependson} \texttt{removedependson} & entre projets \\
\texttt{links} \texttt{removelinks} & bibliothèques \\
\texttt{frameworks} \texttt{framework} & Apple \\
\texttt{defines} \texttt{removedefines} & définitions \\
\texttt{optimize} \texttt{symbols} \texttt{runtime} & réglages de compilation \\
\texttt{cflags} \texttt{cxxflags} \texttt{ldflags} \texttt{arflags} & drapeaux bruts \\
\texttt{prelink} \texttt{postlink} & actions autour du lien \\
\bottomrule
\end{tabular}
\end{center}

\begin{nkretenir}
\textbf{Les \texttt{*flags} sont une porte de sortie, pas un outil courant.} Ils
passent une chaîne au compilateur sans que Jenga la comprenne --- donc sans qu'il
puisse l'adapter à une autre chaîne d'outils.

Un drapeau brut est \textbf{lié à un compilateur}. Employez-le quand rien
d'autre ne convient, dans un filtre \texttt{toolset:}, et écrivez pourquoi.
\end{nkretenir}

\subsection*{Réglages courts}

\begin{center}
\begin{tabular}{@{}ll@{}}
\toprule
\texttt{debug()} \texttt{release()} \texttt{reldebinfo()} \texttt{minsizerel()} & profils \\
\texttt{debuggdb()} \texttt{debuglldb()} \texttt{debugcodeview()} & format de symboles \\
\texttt{sanitizeaddress()} \texttt{sanitizethread()} \dots & analyseurs dynamiques \\
\texttt{coveragegcov()} \texttt{coveragevs()} & couverture \\
\texttt{profilegprof()} \texttt{profilevs()} \texttt{profileinstruments()} & profilage \\
\texttt{pic()} \texttt{pie()} \texttt{nostdlib()} \texttt{nostdinc()} & génération de code \\
\bottomrule
\end{tabular}
\end{center}

\subsection*{Métadonnées et plateformes}

\begin{center}
\begin{tabular}{@{}ll@{}}
\toprule
\texttt{apppublisher} \texttt{appversion} \texttt{licensefile} \texttt{appicon} & communes \\
\texttt{android*} (une trentaine) & Android \\
\texttt{ios*} \texttt{tvos*} \texttt{watchos*} \texttt{visionos*} & Apple \\
\texttt{harmony*} (une quinzaine) & HarmonyOS \\
\texttt{emscripten*} & Web \\
\texttt{xbox*} \texttt{gdkpath} & Xbox \\
\texttt{signing*} & signature \\
\texttt{networkenabled} \texttt{firewallrule} & réseau \\
\bottomrule
\end{tabular}
\end{center}

\begin{nkretenir}
\textbf{Le vocabulaire complet compte 242 fonctions publiques}, relevées dans
\nkfile{Jenga/Core/Api.py}. Les tables ci-dessus en donnent la charpente.

\textbf{Devant un besoin précis, cherchez par préfixe} : \texttt{android},
\texttt{ios}, \texttt{harmony}, \texttt{emscripten}, \texttt{signing}. La
convention de nommage est régulière, et c'est ce qui rend l'API cherchable sans
documentation.
\end{nkretenir}

\section{Les variables de chemin}

\begin{center}
\begin{tabular}{@{}ll@{}}
\toprule
\texttt{\%\{wks.location\}} & le dossier du workspace \\
\texttt{\%\{prj.name\}} & le nom du projet \\
\texttt{\%\{prj.location\}} & le dossier du projet \\
\texttt{\%\{cfg.buildcfg\}} & \texttt{Debug} ou \texttt{Release} \\
\texttt{\%\{cfg.system\}} & le système visé \\
\bottomrule
\end{tabular}
\end{center}

\section{Les préfixes de filtre}

\begin{center}
\begin{tabular}{@{}ll@{}}
\toprule
\texttt{system:} & le système visé \\
\texttt{config:} \texttt{configuration:} \texttt{configurations:} \texttt{cfg:} & la configuration \\
\texttt{arch:} \texttt{architecture:} & l'architecture \\
\texttt{platform:} \texttt{platforms:} & système + architecture \\
\texttt{action:} & la commande en cours \\
\texttt{options:} \texttt{option:} & une option déclarée \\
\texttt{kind:} & le genre du projet \\
\texttt{language:} & le langage \\
\texttt{toolset:} & la famille de compilateur \\
\bottomrule
\end{tabular}
\end{center}

\section{Les onze chaînes d'outils}

\begin{center}
\begin{tabular}{@{}lll@{}}
\toprule
\textbf{Nom} & \textbf{Famille} & \textbf{Cible} \\
\midrule
\texttt{host-apple-clang} & apple-clang & la machine \\
\texttt{host-clang} & clang & la machine \\
\texttt{host-gcc} & gcc & la machine \\
\texttt{clang-mingw} & clang & Windows \\
\texttt{mingw} & gcc & Windows \\
\texttt{msvc} & msvc & Windows \\
\texttt{gcc-cross-linux} & gcc & Linux \\
\texttt{clang-cross-linux} & clang & Linux \\
\texttt{emscripten} & emscripten & Web \\
\texttt{android-ndk} & android-ndk & Android \\
\texttt{ohos-ndk} & clang & HarmonyOS \\
\bottomrule
\end{tabular}
\end{center}

\begin{nkretenir}
\textbf{Déclarées $\neq$ disponibles.} \texttt{jenga info} ne montre que celles
qu'il a trouvées sur la machine. Un descripteur qui exige une chaîne absente
échoue au chargement --- et le message parle de la chaîne, pas du SDK manquant.
\end{nkretenir}

\begin{nkbilan}
Vingt-sept commandes, quinze alias, 242 fonctions de vocabulaire, neuf préfixes
de filtre, onze chaînes d'outils. Les tables de ce chapitre en donnent la
charpente ; le reste se cherche par préfixe, la convention de nommage étant
régulière. Retenez deux choses : les drapeaux bruts sont une porte de sortie liée
à un compilateur, et une chaîne d'outils déclarée n'est pas une chaîne
disponible.
\end{nkbilan}

\begin{nkquestions}
\begin{enumerate}
\item Quelle commande fait \texttt{clean} puis \texttt{build} ?
\item Quels sont les deux alias qui sont des synonymes et non des raccourcis ?
\item Pourquoi un drapeau brut est-il moins portable qu'un réglage nommé ?
\item Que valent \texttt{\%\{cfg.buildcfg\}} et \texttt{\%\{cfg.system\}} ?
\item Quelle différence entre une chaîne d'outils déclarée et disponible ?
\end{enumerate}
\end{nkquestions}

\begin{nkpratique}
Prenez trois besoins et trouvez la fonction, \emph{sans} lire ce chapitre :

\begin{enumerate}
\item embarquer une police dans le binaire ;
\item interdire l'usage de la bibliothèque standard ;
\item déclarer les permissions Android.
\end{enumerate}

\texttt{jenga help}, les exemples fournis et la recherche par préfixe suffisent.
Chronométrez : si l'un vous prend plus de deux minutes, notez lequel --- c'est un
manque de la convention de nommage, et il vaut d'être signalé.
\end{nkpratique}

\begin{nkdemo}
Prouvez qu'une chaîne d'outils déclarée peut être indisponible.

\begin{enumerate}
\item relevez la liste que \texttt{jenga info} affiche sur votre machine ;
\item comparez-la aux onze du tableau ci-dessus ;
\item écrivez un \texttt{usetoolchain} vers une chaîne absente de votre liste et
      lancez \texttt{jenga build}.
\end{enumerate}

\textbf{Ce que la démonstration doit établir} : le message d'échec parle de la
\emph{chaîne}, pas du SDK qui lui manque. Relevez-le mot pour mot --- c'est celui
que vous reverrez sur une machine neuve, et vous saurez alors quoi installer
plutôt que quoi corriger.
\end{nkdemo}

\begin{nklogique}
Vous héritez d'un descripteur qui emploie \texttt{cxxflags(["-O3"])} sous un
filtre \texttt{config:Release}, alors que \texttt{optimize("Speed")} existe.

\begin{enumerate}
\item Citez deux raisons légitimes de ce choix.
\item Citez deux problèmes qu'il crée, dont un qui ne se verra que sur une autre
      plateforme.
\item Que feriez-vous --- et qu'écririez-vous à côté, quel que soit votre choix ?
\end{enumerate}
\end{nklogique}
